\section{Tabela-resumo}
\label{sec:vms-tabela}

A tabela junta as seções anteriores. Cada célula traz a fonte; onde não
encontramos número publicado e comparável, a célula diz isso em vez de
inventar um. Números de desempenho são do fornecedor e servem como ordem
de grandeza (capítulo~\ref{cap:benchmarks}).

\footnotesize
\begin{longtable}{@{}>{\raggedright\arraybackslash}p{1.9cm} >{\raggedright\arraybackslash}p{3.4cm} >{\raggedright\arraybackslash}p{3.4cm} >{\raggedright\arraybackslash}p{3.6cm}@{}}
\toprule
Eixo & HotSpot (JDK~25) & BEAM (OTP~29.1.2) & GraalVM (JIT e Native
  Image) \\
\midrule \endhead
Modelo de concorrência & Threads de plataforma (do SO) e virtual
  threads num \texttt{ForkJoinPool}, sem \emph{time sharing}; heap
  compartilhado (\citetitle{doc:jep444}; \citetitle{doc:jvms25},
  §2.5.3) & Processos leves, sem memória compartilhada, mensagens
  copiadas para a mailbox; preempção por reductions
  (\citetitle{doc:erlang-processes};
  \citetitle{doc:erlang-eff-processes}; \citetitle{livro:stenman}) &
  O da HotSpot, na qual a GraalVM se baseia
  (\citetitle{doc:graalvm-java}); no Native Image, as mesmas APIs Java
  \\
\addlinespace
GC & Um heap para todas as threads; G1 padrão em máquinas
  server-class (Serial nas outras), ZGC, Parallel
  (seção~\ref{sec:gc}; \citetitle{doc:gc-collectors}) & Por processo,
  geracional por cópia (Cheney), mais espaço global de binários
  contados por referência (\citetitle{doc:erlang-gc}) & JIT: os da
  HotSpot. Native: Serial (padrão, \emph{stop and copy}); G1 só no
  Oracle GraalVM em Linux; Epsilon (\citetitle{doc:graalvm-memory}) \\
\addlinespace
Tratamento de falha & Exceção não tratada termina a thread; o estado no
  heap fica (\citetitle{doc:jls25}, §11.3); resiliência por bibliotecas
  (seção~\ref{sec:resiliencia-quarkus}) & \emph{Let it crash}: links,
  sinais de saída e árvores de supervisão com limite de reinícios
  (\citetitle{paper:armstrong-thesis-reliable-systems};
  \citetitle{doc:erlang-supervisor}) & Igual à HotSpot nos dois modos:
  é a mesma linguagem Java (\citetitle{doc:jls25}, §11.3;
  \citetitle{doc:graalvm-java}) \\
\addlinespace
Carga de código & Class loading dinâmico
  (seção~\ref{sec:class-loading}); redefinição só de corpos de métodos
  (\citetitle{doc:jvmti25}) & Duas versões por módulo; chamada
  qualificada migra para a nova; a terceira carga mata quem ficou na
  mais antiga (\citetitle{doc:erlang-code-loading}) & JIT: como a
  HotSpot. Native: mundo fechado, nada novo em runtime
  (\citetitle{doc:graalvm-native-basics}) \\
\addlinespace
Compilação & Interpretador + C1/C2 com perfil, especulação e
  desotimização (seção~\ref{sec:jit}; \citetitle{doc:hotspot-perf}) &
  BeamAsm: tradução na carga, sem otimização entre instruções
  (\citetitle{doc:erlang-beamasm}) & JIT: Graal, escrito em Java, com
  \emph{partial escape analysis} (\citetitle{doc:graalvm-compiler}).
  Native: AOT após análise de \emph{points-to}
  (\citetitle{paper:native-image-initialize-once}) \\
\addlinespace
Startup e memória & Quarkus em modo JVM: 4417~ms até a 1ª requisição,
  304~MiB de RSS. Perf-lab do Quarkus, 21/04/2026, Quarkus~3.34.3,
  JDK~25.0.2 (fornecedor; seção~\ref{sec:jvm-vs-nativo}) & 327
  palavras por processo novo (OTP~27, \citetitle{doc:erlang-eff-processes}).
  Não achamos medição publicada de startup e RSS de um serviço
  comparável ao do Quarkus & Native: 581~ms e 95~MiB na mesma medição
  do Quarkus (fornecedor; seção~\ref{sec:jvm-vs-nativo}). O paper do
  Native Image fala em startup \enquote{up to two orders of magnitude}
  melhor que a HotSpot (\citetitle{paper:native-image-initialize-once};
  Oracle Labs, fornecedor, 2019) \\
\bottomrule
\end{longtable}
\normalsize

\paragraph{Como ler a tabela.}
A linha de startup é a única com números, e eles não se comparam entre
colunas: a medição do Quarkus é de um serviço REST com banco, e as 327
palavras da BEAM são o custo de \emph{um} processo, não de um serviço.
O que a tabela mostra bem são as escolhas de projeto: a HotSpot aposta em
heap compartilhado e JIT especulativo; a BEAM, em isolamento e troca a
quente; o Native Image, em saber tudo no build.
